\section{Казус: непривързаните литерални операнди}
\label{sec:case-study}

Глава~\ref{sec:enforcement} установи, че всеки конектор свързва стойностите на
заместителите правилно. Тази глава проследява случай, в който предпоставка 2 от
точка~\ref{sec:theory-compile-time} не е била изпълнена, въпреки че всеки
конектор поотделно работи коректно. Казусът е избран по две причини. Той
демонстрира метода на работата върху реален дефект, от хипотезата през
измерването до поправката. Той показва и защо анализът на всеки компонент
поотделно не е достатъчен, когато гаранцията зависи от това, какво изобщо може да
достигне компонентите.

\subsection{Откриване}
\label{sec:cs-discovery}

Отправната точка беше въпросът дали типизираният заместител е
\textit{единственият} път, по който стойност достига заявката. Операторите за
сравнение на полета са дефинирани чрез макрос в слоя на заявките.
Листинг~\ref{lst:cs-operator} показва съществената му част.

\begin{lstlisting}[language=C++,caption={Операторът за сравнение на поле (съкратено)},label={lst:cs-operator}]
template <auto P, typename T>
    requires(!std::is_same_v<T, std::nullptr_t>)
constexpr auto operator==(mem_ptr<P>, T x)
{
    static_assert(comparable_value_types<...>, "...");
    if constexpr (is_field<T>)
        return Rule<decltype(P), "==", decltype(T::get())>(P, T::get());
    else
        // стойността се записва в заявката
        return Rule<decltype(P), "==", T>(P, x);
}
\end{lstlisting}

От листинга се вижда, че параметърът \code{T} не е ограничен до заместители.
Операторът приема всяка стойност със съвместим тип и я записва в междинното
представяне като десен операнд. Следователно съществува и втори път: стойност,
записана направо в условието, например \code{field<\&Item::name> == input}, без
заместител. Такава стойност се нарича в тази работа \textbf{литерален операнд}.
Ако тя бъде вмъкната в текста на заявката, гаранцията е нарушена. Ако бъде
свързана, гаранцията е запазена. Кодът на оператора не дава отговор, тъй като
решението се взема в конектора.

\subsection{Хипотеза и измерване}
\label{sec:cs-measurement}

Конекторът за SQLite превежда всеки операнд, който не е колона и не е
заместител, в маркер \code{?} (листинг~\ref{lst:cs-render}, последният клон).
Литералът не става текст, което изключва инжекцията. Търсене на код, който след това свързва
стойността на литерала, не откри такъв в нито един от осемте конектора (седемте
работещи и \code{MockDB}). Хипотезата беше, че маркерът остава несвързан и
драйверът го тълкува като \code{NULL}.

\begin{lstlisting}[language=C++,caption={Преводът на операнд при SQLite преди поправката (съкратено)},label={lst:cs-render}]
template <typename T>
std::string render_operand_sqlite(const T& /*v*/)
{
    if constexpr (is_field<T>)
        return std::string(T::column_name());
    // ...
    else if constexpr (is_placeholder_v<T>)
    {
        // ...
    }
    else
        return "?";   // литерален операнд: маркер без стойност
}
\end{lstlisting}

Хипотезата беше проверена с две програми върху SQLite в паметта, с библиотеката
във версията преди поправката (\code{d5795a9}). Таблица~\ref{tab:cs-select}
показва резултата за заявка за избор.

\begin{table}[H]
\centering
\caption{Заявка за избор: заместител спрямо литерал}
\label{tab:cs-select}
\begin{tabularx}{\textwidth}{L{4.2cm}L{4.2cm}L{2.2cm}Y}
\toprule
\textbf{Път} & \textbf{Вход} & \textbf{Редове} & \textbf{Правилно} \\
\midrule
типизиран заместител & \code{banana} & 1 & 1 \\
типизиран заместител & \code{x' OR '1'='1} & 0 & 0 \\
литерал в заявката & \code{banana} & \textbf{0} & 1 \\
литерал в заявката & \code{x' OR '1'='1} & 0 & 0 \\
\bottomrule
\end{tabularx}
\end{table}

Резултатът потвърждава хипотезата. Пътят през заместителя връща правилния
отговор и при безобиден вход, и при полезен товар. Пътят през литерала връща
нула редове дори за \code{banana}, за който съществува съвпадащ ред. Полезният
товар не постига нищо и по двата пътя, но това не означава едно и също: при
литерала неуспехът му се дължи на това, че заявката изобщо не работи.

Таблица~\ref{tab:cs-update} показва същото при промяна и изтриване.

\begin{table}[H]
\centering
\caption{Промяна и изтриване: заместител спрямо литерал}
\label{tab:cs-update}
\begin{tabularx}{\textwidth}{L{5.4cm}L{4.4cm}Y}
\toprule
\textbf{Случай} & \textbf{Намерение} & \textbf{Резултат} \\
\midrule
заместител в \code{SET} & преименуване на ред 1 & редът е преименуван \\
литерал в \code{SET} & преименуване на ред 2 & \textbf{редът е непроменен} \\
литерал в \code{WHERE} & изтриване на ред 3 & \textbf{нищо не е изтрито} \\
\bottomrule
\end{tabularx}
\end{table}

Листинг~\ref{lst:cs-probe} показва дословния изход на втората програма.

\begin{lstlisting}[language={},deletekeywords={orm},caption={Изход на програмата за промяна и изтриване},label={lst:cs-probe}]
U1 set via placeholder 'grape', id=1   -> name now grape
U2 set via literal     'grape', id=2   -> name now banana
D1 delete via literal id==3            -> rows left 3 (3 means nothing deleted)
\end{lstlisting}

\subsection{Анализ на причината}
\label{sec:cs-cause}

Случаят с промяната заслужава отделно внимание, защото резултатът му не е
просто липса на действие. Аргументите, подадени на \code{execute()}, се свързват
позиционно, от позиция 1 нататък (листинг~\ref{lst:enf-sqlite-bind}). Заявката
\code{UPDATE items SET name = ? WHERE id = ?} има два маркера, а при изпълнението
е подаден един аргумент, предназначен за условието. Маркерът на литерала в
\code{SET} е първи. Затова аргументът е свързан с него, а маркерът в условието е
останал несвързан, тоест \code{NULL}, и не е съвпаднал с нито един ред. Това е
\textbf{объркване на параметрите}: стойност на извикващия код попада на погрешна
позиция.

От гледна точка на сигурността дефектът е \textbf{затворен отказ}: при грешка
системата остава в безопасно състояние. Литералът никога не става текст и не
може да промени структурата. Дефектът обаче не е безобиден. Той води до тихо
грешни резултати, без грешка и без предупреждение. За приложение, което
проверява правомощия чрез заявка, грешен празен резултат може да означава отказ
на легитимен достъп. Случай с видим погрешен ефект, например промяна на погрешен
ред, не беше наблюдаван в измерените заявки, но не е изключен за всяка форма на
заявка.

Остава въпросът защо дефектът не е бил открит от съществуващите тестове. Търсене
в тестовете намери литерални операнди само в единичните тестове на междинното
представяне, които проверяват формата на междинното представяне чрез
\code{static\_assert} и никога не изпълняват заявката. Нито един тест не е
изпълнявал заявка с литерален операнд срещу база данни. Дефектът е оцелял в
сляпо петно на тестовото покритие, а не поради проектно решение.

\subsection{Избор на поправка}
\label{sec:cs-fix}

Възможни са две поправки.

\textbf{Свързване на литералите.} Всеки конектор събира стойностите на
литералите и ги свързва заедно с аргументите на \code{execute()}, в общ ред на
позициите. Така литералите стават законен втори път за стойности.

\textbf{Отхвърляне на литералите.} Заявка с литерален операнд не се компилира, ако
бъде подадена за изпълнение. Единственият път за стойности остава заместителят.

Изборът падна върху отхвърлянето по три причини. Първо, свързването изисква
промяна в превода и свързването на всички седем конектора. Грешка в реда на
позициите в който и да е от тях би превърнала затворения отказ в объркване на
параметрите с видим ефект, а реализация, която вмъква литерала като текст, би
създала истинска инжекция. Отхвърлянето е една проверка на едно място
(точка~\ref{sec:enf-boundary}). Второ, то съответства на тезата на работата:
стойностите влизат само през заместители. Трето, поправката не чупи работещ код,
защото кодът, който тя отхвърля, и без това е давал грешни резултати.

Цената е, че естествени записи, например условие за положително число, вече
изискват заместител. Отхвърлянето засяга и литерали, известни по време на
компилация, защото дефектът никога не е зависел от това кога е известна
стойността: нито един конектор не свързва литерал.

\subsection{Проверка на поправката}
\label{sec:cs-verification}

Поправката е проверена в два етапа. Преди нея наборът от компилационни тестове
беше изпълнен с библиотеката във версия \code{d5795a9}. Седемте теста за
литерални операнди, по един за всяка входна точка и вид заявка, се провалиха с
един и същ резултат: програмата се компилира, въпреки че не би трябвало.
От 21 теста в това изпълнение 12 преминаха, а 9 се провалиха: седемте за
литерални операнди и двата, разгледани в точка~\ref{sec:eval-compile}.

След поправката и седемте програми се отхвърлят при компилация по проверяваната
причина. Всички тестове преминават: 314-те съществуващи в библиотеката и 23-те
нови от набора за оценка. Последното е съществено: то показва, че поправката
затваря дефекта, без да засегне работещ код. Подробните резултати са в
глава~\ref{sec:evaluation}.

\subsection{Изводи от казуса}
\label{sec:cs-lessons}

Казусът води до два извода, които надхвърлят конкретния дефект. Първият е, че
гаранцията за сигурност на една абстракция не е сума от гаранциите на
компонентите ѝ. Всеки конектор свързва заместителите правилно, но никой от тях не
е отговарял за въпроса какво изобщо може да ги достигне. Вторият е, че тест,
който проверява само формата на заявката, не проверява поведението ѝ. Дефектът
оцеля, защото литералите бяха тествани като структура, но не и като изпълнени
заявки. Методът на настоящата работа, който изисква всяко твърдение да бъде
измерено, а не изведено, е причината дефектът да бъде открит.
